\documentclass[12pt]{article}
\usepackage[utf8]{inputenc}
\usepackage{amsmath, amssymb, amsthm}
\usepackage{mdframed}
\usepackage{geometry}

\geometry{a4paper, margin=1in}

\newtheorem{definition}{Definition}

\parskip=1ex
\title{Proofs with Sets}
\date{}

\begin{document}

\maketitle

$A = \{0,1,2,3,4,5\}$ and $R = \{ (1,2), (3,3), (2,5), (4,2), (4,3), (5,0) \}$

\begin{definition}
    A relation $R$ on $\mathbf{X}$ is \textbf{reflexive} if for all $a \in \mathbf{X}, aRa$
\end{definition}
\begin{definition}
    A relation $R$ on $\mathbf{X}$ is \textbf{symmetric} if for all $a,b \in \mathbf{X}$, if $aRb$ then $bRa$.
\end{definition}
\begin{definition}
    A relation $R$ on $\mathbf{X}$ is \textbf{transitive} if for all $a,b,c \in \mathbf{X}$, if $aRb$ and $bRc$, then $aRc$.
\end{definition}


Is $R$ reflexive? Explain.
\begin{proof}
    Say we choose $m=2$.

    Since $gcd(2, 2) = 2 \neq 1$, it follows that $mRm$ is false.

    Thus $R$ is not reflexive
\end{proof}

Is $R$ symmetric? Explain.
\begin{proof}
    Choose arbitrary natural numbers $a$ and $b$.

    Suppose $aRb$.

    Then, we have that $gcd(a, b) = 1$.

    Thus, $gcd(b, a ) = 1$.

    So $bRa$.

    Therefore, we conclude that $R$ is symmetric.
\end{proof}

Is $R$ transitive? Explain.
\begin{proof}
    Consider $2,3,4 \in \mathbb{N}$.

    Since $gcd(2,3) = 1$, it follows that $2R3$.
    Since $gcd(3,4) = 1$, it follows that $3R$.

    But, $gcd(2,4) = 2 \neq 1$. So $2\not\mathrel{R}$.
\end{proof}

\end{document}